\chapter{Dataset and Task Definition}
\label{chapter:dataset}

This chapter describes the dataset used throughout the thesis and defines the classification
tasks evaluated in the experimental chapters. The dataset consists of simulated
proton-proton collision events at $\sqrt{s} = 13\,\text{TeV}$, covering one signal process
($t\bar{t}t\bar{t}$, hereafter four-tops) and four Standard Model (SM) background processes
($t\bar{t}H$, $t\bar{t}W$, $t\bar{t}WW$, $t\bar{t}Z$). Events are encoded as fixed-length
token sequences and stored in HDF5 format. The dataset was received pre-split and used
without modification to the train/validation/test partition. All architectural experiments
in later chapters operate on this single dataset and evaluate binary signal-versus-background
discrimination.

% -------------------------------------------------------------------------
\section{Dataset Provenance}
\label{sec:dataset_provenance}
% -------------------------------------------------------------------------

The dataset originates from two upstream sources. The event generation and detector
simulation pipeline is documented by Builtjes et al.\
\cite{builtjesAttentionStrengthsPhysical2025}, who produced the dataset to study
transformer and graph-based classifiers on multi-top final states. The data format follows
the convention established by the Dark Machines collaboration \cite{Aarrestad:2021oeb}, a
community effort that defined a standardised benchmark format for model-independent LHC
event classification.

Hard-scattering events were generated at leading order using \textsc{MadGraph5\_aMC@NLO}
2.7 with the NNPDF3.1 parton distribution function set and the five-flavour scheme. Parton
showering and hadronisation were performed with \textsc{Pythia}~8.239 using the MLM merging
scheme. Fast detector simulation was carried out with \textsc{Delphes}~3.4.2 using the
ATLAS detector card. An event-generation-level requirement of $H_T > 400\,\text{GeV}$,
summed over parton-level jets, was imposed to increase simulation efficiency for processes
whose reconstructed signal region demands $H_T > 500\,\text{GeV}$. Generated Monte Carlo
samples were converted from ROOT to the Dark Machines CSV format, from which the dataset
used in this thesis was constructed.

The dataset was received in pre-split HDF5 format (\texttt{4tops\_splitted.h5}), with no
train/validation/test partition performed in this work. The exact splitting procedure,
including stratification criterion and random seed, is not documented in the source
materials. \todo{ask supervisor: splitting methodology}

% -------------------------------------------------------------------------
\section{Data Format and Representation}
\label{sec:dataset_format}
% -------------------------------------------------------------------------

Each event in the original CSV follows the Dark Machines line format
\cite{Aarrestad:2021oeb,builtjesAttentionStrengthsPhysical2025}:
\begin{center}
event ID;\enspace process ID;\enspace weight;\enspace
$\|\vec{E}_T^{\,\text{miss}}\|$;\enspace $\phi_{E_T^{\text{miss}}}$;\enspace
$\mathrm{obj}_1,\,E_1,\,p_{T,1},\,\eta_1,\,\phi_1$;\enspace
$\mathrm{obj}_2,\,E_2,\,p_{T,2},\,\eta_2,\,\phi_2$;\enspace$\ldots$
\end{center}
Each reconstructed particle object is characterised by an integer type tag and its
four-momentum $(E,\,p_T,\,\eta,\,\phi)$. Seven object types are defined: jets (\texttt{j},
tag~1), $b$-tagged jets (\texttt{b}, tag~2), positrons (\texttt{e+}, tag~3), electrons
(\texttt{e-}, tag~4), muons (\texttt{mu-}, tag~5), anti-muons (\texttt{mu+}, tag~6), and
photons (\texttt{g}, tag~7). The missing transverse energy magnitude
$\|\vec{E}_T^{\,\text{miss}}\|$ and its azimuthal angle $\phi_{E_T^{\text{miss}}}$ are
stored as two global features shared across all objects in the event.

Because the number of reconstructed objects varies per event, all sequences are
zero-padded to a fixed length of $T = 18$ token slots \cite{Visive:2026mjd}. The value
$T = 18$ corresponds to the largest event in the dataset. A binary attention mask
accompanies each event: padding tokens receive a mask value of zero and are excluded from
attention computations, while real object tokens receive a mask value of one.

The HDF5 file stores each event as a flat row of 92 columns: 18 integer type IDs,
2 global features ($\|\vec{E}_T^{\,\text{miss}}\|$ and $\phi_{E_T^{\text{miss}}}$), and
$18 \times 4 = 72$ continuous kinematic features (one four-momentum per token slot).
The integer type IDs and continuous features are separated at runtime by the data loader.
Before training, the continuous features are normalised using global z-score statistics
computed over the training split only, pooled across all $N_{\text{train}}$ events and all
$T$ token positions simultaneously. This yields one scalar mean and one scalar standard
deviation per feature type ($E$, $p_T$, $\eta$, $\phi$), applied uniformly to every token
slot in every event. No per-position or per-particle-type rescaling is performed.

Across 98.6\% of all thesis training runs, tokens are presented in the order they appear
in the HDF5 file (\texttt{input\_order}). No reordering is applied by the data loader for
these runs. The provenance of the within-event object ordering in the source file is
unknown. \todo{ask supervisor: within-event object ordering in HDF5}

A discrete tokenisation scheme, referred to as AmbreBinning and developed for a subset of
experiments in Chapter~4, maps each continuous four-momentum to a single integer through
independent uniform binning of $p_T$, $\eta$, and $\phi$ (five bins each) combined with
the object-type dimension (125 categories). The combined token index is computed as
$125\,(b_{\text{obj}}-1) + 25\,(b_{p_T}-1) + 5\,(b_\eta-1) + b_\phi$, yielding a
vocabulary of 886 distinct integers. This scheme was inspired by binning strategies
introduced by the thesis supervisor for a related dataset and is not the default
representation; the majority of experiments use continuous four-momentum features directly.

% -------------------------------------------------------------------------
\section{Classification Tasks}
\label{sec:dataset_tasks}
% -------------------------------------------------------------------------

The dataset contains 302,072 events distributed across five SM process classes, split
80/10/10 into training, validation, and test subsets.
Table~\ref{tab:sample_counts} gives the per-class event counts for each split.

\begin{table}[h]
\centering
\caption{Event counts per process class and dataset split. Label codes follow the
convention of Builtjes et al.\ \cite{builtjesAttentionStrengthsPhysical2025}:
1\,=\,$t\bar{t}t\bar{t}$, 2\,=\,$t\bar{t}H$, 3\,=\,$t\bar{t}W$,
4\,=\,$t\bar{t}WW$, 5\,=\,$t\bar{t}Z$.}
\label{tab:sample_counts}
\begin{tabular}{lccccc|c}
\hline
\textbf{Split} & $t\bar{t}t\bar{t}$ & $t\bar{t}H$ & $t\bar{t}W$ & $t\bar{t}WW$ &
$t\bar{t}Z$ & \textbf{Total} \\
\hline
Train & 121{,}359 & 28{,}854 & 30{,}511 & 30{,}461 & 30{,}473 & 241{,}658 \\
Val   & 15{,}266  & 3{,}666  & 3{,}779  & 3{,}760  & 3{,}736  & 30{,}207  \\
Test  & 15{,}375  & 3{,}552  & 3{,}710  & 3{,}779  & 3{,}791  & 30{,}207  \\
\hline
\end{tabular}
\end{table}

The signal process, $t\bar{t}t\bar{t}$, constitutes approximately 50\% of all events. The
four background processes ($t\bar{t}H$, $t\bar{t}W$, $t\bar{t}WW$, $t\bar{t}Z$) each
contribute roughly 10--12\% of the total. This deliberate 1:1 signal-to-background ratio,
chosen by Builtjes et al.\ to balance training \cite{builtjesAttentionStrengthsPhysical2025},
differs markedly from the physical production cross-sections, which strongly favour the
backgrounds, and should be kept in mind when interpreting classifier outputs.

The primary classification task evaluated throughout this thesis is binary discrimination
of $t\bar{t}t\bar{t}$ (label 1) against the combined SM background (labels 2--5). This is
the most physically relevant task: the four-tops process has an extremely small production
cross-section ($\sigma \approx 10\,\text{fb}$ at $\sqrt{s} = 13\,\text{TeV}$), and its
first observation by ATLAS required careful discrimination from a large and topologically
similar SM background \cite{collaborationObservationFourtopquarkProduction2023}. The task
is challenging because each top quark decays almost exclusively via $t \to Wb$, so four-top
events produce high-multiplicity final states with 0--4 leptons and 4--12 jets, including
four $b$-jets. Background processes such as $t\bar{t}WW$ and $t\bar{t}Z$ produce
comparably complex topologies. The classifier must therefore exploit subtle differences in
object multiplicity, kinematic correlations, and $b$-jet content to separate signal from
background.

All architectural ablations in Chapters~4--8 evaluate this binary task, reporting area
under the receiver operating characteristic curve (AUROC) and background rejection at fixed
signal efficiency as primary metrics.
